\documentclass[../DoAn.tex]{subfiles}
\begin{document}

\section{Đặt vấn đề}
\label{section:1.1}

Đến năm 2026, số mẫu mã độc trên Windows do AV-TEST ghi nhận đã vượt
1 tỷ và tiếp tục tăng đều qua từng năm~\cite{avtest2025}. Theo báo cáo IBM~\cite{ibm2025}, chi phí trung
bình mỗi vụ vi phạm dữ liệu đạt 4,44 triệu USD năm 2025, bao gồm tổn
thất do gián đoạn hoạt động, mất khách hàng và chi phí tuân thủ pháp lý.
Ở quy mô toàn cầu, thiệt hại do tội phạm mạng được dự báo đạt 10,5
nghìn tỷ USD mỗi năm~\cite{cybercrime2025}.

Phần mềm độc hại trên Windows tồn tại dưới nhiều hình thức với cơ chế
tấn công khác nhau. Ransomware mã hóa dữ liệu rồi đòi tiền chuộc và
hiện diện trong 48\% số vụ vi phạm dữ liệu theo báo cáo Verizon DBIR
2026~\cite{verizon2026}. Trojan ngụy trang thành phần mềm hợp lệ để cài
cửa hậu hoặc đánh cắp thông tin; biến thể phổ biến là RAT, duy trì
quyền điều khiển từ xa trung bình 11 ngày trước khi bị phát
hiện~\cite{mandiant2025}. Worm tự nhân bản và lan truyền qua mạng mà
không cần thao tác của người dùng nên phát tán rất nhanh: mã độc dạng
worm như WannaCry từng lây nhiễm hơn 200.000 máy tính tại ít nhất 150
quốc gia chỉ trong vài ngày, gây thiệt hại ước tính tới 4 tỷ
USD~\cite{wannacry2017}. Sự đa dạng về cơ chế tấn công đòi hỏi phương
pháp phát hiện thích ứng được với nhiều dạng mối đe dọa thay vì chỉ
dựa vào một cơ chế duy nhất.

Phương pháp phát hiện phổ biến nhất hiện nay là nhận diện dựa trên chữ
ký: hệ thống đối chiếu tệp cần kiểm tra với cơ sở dữ liệu mẫu độc hại
đã biết. Cách làm này hiệu quả với các mối đe dọa đã biết nhưng bộc lộ
hạn chế trước biến thể mới, bởi kẻ tấn công chỉ cần đóng gói lại hoặc
làm rối mã thực thi là đủ khiến chữ ký hiện có không còn nhận diện
được~\cite{gibert2020}. Riêng năm 2024, SonicWall ghi nhận hơn 210.000
biến thể mã độc mới~\cite{sonicwall2025}; con số này phản ánh khoảng
trễ thường trực giữa thời điểm mẫu xuất hiện và thời điểm chữ ký được
cập nhật.

Hạn chế này đặt ra nhu cầu về phương pháp phát hiện không phụ thuộc
hoàn toàn vào chữ ký, có khả năng nhận diện cả biến thể zero-day. Kết
hợp phân tích tĩnh bằng học máy và quan sát hành
vi động trong môi trường sandbox là hướng tiếp cận cho phép hệ thống
nhận diện cả mẫu đã biết lẫn biến thể mới.

\section{Mục tiêu và phạm vi đề tài}
\label{section:1.2}

Từ hạn chế của phương pháp phát hiện dựa trên chữ ký đã phân tích tại
Mục~\ref{section:1.1}, đồ án xây dựng hệ thống hỗ trợ phát hiện và
phân tích phần mềm độc hại trên nền tảng Windows, thực hiện phân loại
tệp PE thành lành tính hoặc độc hại đồng thời thu thập bằng chứng hành
vi phục vụ điều tra.
Cụ thể, bên cạnh kết quả phân loại, hệ thống thu thập thông tin phục
vụ phân tích hành vi gồm IOC, cây tiến trình và các
kỹ thuật tấn công theo khung MITRE ATT\&CK~\cite{mitre_attack}. Ba dạng đầu ra này bổ sung
cho nhau: IOC phục vụ truy vết và đối chiếu nhanh, cây tiến trình giúp
hiểu luồng thực thi, ánh xạ ATT\&CK chuẩn hóa kết quả cho đội ngũ
phân tích.

Đồ án giới hạn phạm vi đối với các tệp thực thi định dạng PE 32/64-bit
(\texttt{.exe}, \texttt{.dll}) trên hệ điều hành Windows. Các định dạng
khác như tài liệu Office chứa macro, script PowerShell hay JavaScript,
cũng như các nền tảng ngoài Windows, không thuộc phạm vi đề tài. Đồ
án không đặt mục tiêu xử lý các kỹ thuật né tránh sandbox nâng cao
hay cơ chế chống phân tích chuyên biệt.

\section{Định hướng giải pháp}
\label{section:1.3}

Để đạt được các mục tiêu đề ra, đồ án đề xuất kiến trúc hai giai đoạn
kết hợp phân tích tĩnh và phân tích động. Giai đoạn tĩnh thực hiện
phân loại tệp PE thông qua mô hình học máy huấn luyện trên các đặc
trưng cấu trúc. Giai đoạn động thực thi mẫu trong môi trường sandbox
cô lập nhằm thu thập thông tin hành vi thực tế. Phân tích tĩnh đơn lẻ
dễ bị vô hiệu hóa bởi kỹ thuật đóng gói và làm rối mã; phân tích động
đơn lẻ tiêu tốn nhiều tài nguyên và không cung cấp phân loại tức thời.
Kết hợp hai giai đoạn cho phép hệ thống phân loại nhanh dựa trên cấu
trúc tệp, đồng thời thu thập bằng chứng hành vi để làm rõ ngữ cảnh
tấn công. So với các hệ thống phân tích động hiện có, đồ án tích hợp
thêm SilkETW để giám sát .NET Runtime — lớp hành vi mà Sysmon không
ghi nhận trực tiếp — và hợp nhất kết quả hai giai đoạn thành một báo
cáo thống nhất gồm kết luận, IOC, cây tiến trình và ánh xạ MITRE
ATT\&CK.

\section{Bố cục đồ án}
\label{section:1.4}

Phần còn lại của đồ án được tổ chức như sau.
Chương~\ref{chapter:Related_works} trình bày cơ sở lý thuyết, bao gồm
đặc điểm các nhóm mã độc phổ biến, cấu trúc tệp PE và đặc trưng tĩnh,
cơ chế thu thập sự kiện qua Sysmon và ETW, và khung MITRE ATT\&CK.
Chương~\ref{chapter:Design} trình bày kiến trúc tổng thể, quy trình xây
dựng mô hình học máy, thiết kế môi trường sandbox và cơ chế xây dựng
chuỗi hành vi.
Chương~\ref{chapter:Implementation} trình bày triển khai từng thành phần
và giao diện web demo.
Chương~\ref{chapter:Experiment} đánh giá hiệu năng mô hình phân loại
và kết quả phân tích hành vi động.
Phần Kết luận tóm tắt kết quả đạt được và đề xuất
hướng phát triển.

\end{document}
